\subsection{Short Rollouts and Ensembles: The Modern Dyna}
\label{section:fc_mbpo}

The line of work culminating in \citet{janner2019model} is the continuous-control inheritor of Dyna (\S\ref{section:fc_dyna_q}). The architectural recipe is a deep ensemble of probabilistic dynamics models paired with a model-free off-policy actor-critic, soft actor-critic in the canonical implementation, with synthetic experience generated by branching short imagined rollouts from real states drawn out of the replay buffer. The framing slogan, present in the paper and in the surrounding literature, is to not trust the model too far. Breaking the dependence between the task horizon and the model rollout horizon is the move that makes the model-error penalty tractable, and is what separates this line from the trajectory-optimization approach of \citet{Chua2018PETS} and the on-policy approach of \citet{Kurutach2018}. The conceptual lineage runs back to \citet{sutton1990}, with the same outer loop of direct learning from real transitions and indirect learning from imagined ones, only with a deep ensemble in place of the tabular model and an off-policy actor-critic in place of tabular Q-learning.

\subsubsection{Branched short rollouts}
\label{section:fc_mbpo_branching}

As an abstract operator, the branched rollout maps the pair (real-state distribution, learned dynamics) to a synthetic transition distribution of length $k$, and decouples the task horizon from the model horizon by anchoring every imagined trajectory at a state the agent has actually visited. The MBPO branching scheme is straightforward. At each gradient step the algorithm samples a real state $s$ from the replay buffer, rolls the learned dynamics ensemble forward for $k$ steps from $s$ under the current policy, and trains the actor-critic on the resulting model-generated transitions alongside real ones. The actor-critic update does not distinguish synthetic transitions from real transitions; the only architectural change relative to model-free soft actor-critic is the distribution from which training tuples are drawn. Real states act as anchors. Because every imagined trajectory is no more than $k$ steps long and originates at a state the agent has actually visited, the model is never asked to extrapolate from its own predictions for more than $k$ steps, even when the task itself runs for hundreds of steps.

The returns bound in \citet[Theorem 4.2]{janner2019model} formalizes the trade-off. With $\epsilon_m$ the model error measured in total variation, $\epsilon_\pi$ the policy distribution shift relative to the data-collection policy, $k$ the branching length, and $r_{\max}$ the reward bound,
\[ \eta[\pi] \ge \eta^{\text{branch}}[\pi] - 2 r_{\max} \Big[ \frac{\gamma^{k+1} \epsilon_\pi}{(1 - \gamma)^2} + \frac{\gamma^k + 2}{1 - \gamma} \epsilon_\pi + \frac{k}{1 - \gamma}(\epsilon_m + 2 \epsilon_\pi) \Big]. \]
The first two terms penalize policy shift and shrink in $k$; the third penalizes model error and grows linearly in $k$. Theorem 4.3 then shows that an optimal $k^\star > 0$ exists whenever the model error is small enough, justifying short rollouts as the right way to spend a fixed model-error budget. The practical takeaway, stated in the paper's ablations, is that $k = 1$ is competitive across the MuJoCo suite, that $k$ in the range of roughly five to fifteen typically gives the best returns, and that values much beyond about $25$ are too inaccurate to help. The empirical sweet spot is well inside the regime where the model is treated as a local data augmentation rather than as a long-horizon simulator.

Read in economics terms, the branched-rollout architecture is the computational version of using a locally valid structural model for counterfactual policy search; the model is trusted only over short windows around states the agent has already visited, and the rest of the long-horizon credit is left to the off-policy critic. The dynamics ensemble used in this line, typically five to seven neural networks each predicting a Gaussian over the next-state delta, plays two roles. As an aggregate predictor the ensemble averages over members and trims tails to give a calibrated forecast that a single deterministic network would not produce; as a disagreement signal the spread across members is a cheap proxy for epistemic uncertainty, large where the agent has not visited often and small where coverage is good.\footnote{The ensemble appears in two prior methods with different commitments. \citet{Kurutach2018} pairs the ensemble with on-policy trust-region policy optimization, estimating value gradients while clipping updates that would push the data distribution into regions where members disagree; the on-policy constraint bounds $\epsilon_\pi$ explicitly but keeps the long-horizon dependence on the model. \citet{Chua2018PETS} wraps the same ensemble in model-predictive control and bypasses policy learning entirely by propagating particle trajectories through ensemble members at planning time, which solves the horizon problem by conflating the task and planning horizons and limits scaling to long-horizon tasks where MPC becomes prohibitive. MBPO's branched-rollout architecture pushes the two terms of the returns bound into a regime where standard off-policy machinery can absorb them, and the same recipe carries from low-dimensional locomotion to humanoid without changes to the outer loop.}\footnote{On the MuJoCo continuous-control benchmark \citet{janner2019model} matches the asymptotic return of soft actor-critic with roughly an order of magnitude fewer environment samples; the Ant task reaches the soft actor-critic return at around $300{,}000$ real steps against the latter's roughly $3{,}000{,}000$, and on Humanoid (a $376$-dimensional state with a $17$-dimensional action) MBPO converges where PETS fails to learn, consistent with the architectural reading that the task horizon dwarfs PETS's planning horizon. \citet{Kurutach2018} reports a similar order-of-magnitude sample-efficiency gain for ME-TRPO over TRPO/PPO/DDPG on the same suite.}
